% Generated draft: requires mathematical and source review.
\documentclass[11pt]{article}
\usepackage[margin=25mm]{geometry}
\usepackage{fontspec}
\setmainfont{Latin Modern Roman}
\usepackage{amsmath,amssymb,mathtools,unicode-math}
\setmathfont{Latin Modern Math}
\usepackage{enumitem,needspace,xurl,hyperref}
\setlength{\parindent}{0pt}
\setlength{\parskip}{5pt}
\setlength{\emergencystretch}{3em}
\newcommand{\N}{\mathbb{N}}
\newcommand{\Z}{\mathbb{Z}}
\newcommand{\Q}{\mathbb{Q}}
\newcommand{\R}{\mathbb{R}}
\newcommand{\C}{\mathbb{C}}
\newcommand{\F}{\mathbb{F}}
\newcommand{\A}{\mathbb{A}}
\newcommand{\PP}{\mathbb{P}}
\newcommand{\E}{\mathbb{E}}
\newcommand{\eps}{\varepsilon}
\newcommand{\dd}{\,\mathrm{d}}
\begin{document}
\section*{UnsolvedMath ID: 800533}
\textbf{Source problem code:} catalogue entry
\subsection{Definitions and mathematical statement}
Let \(G=(V,E)\) be a finite simple undirected graph, supplied by its adjacency matrix. For distinct points \(p,q\in\R^2\), define their open segment by
\[
(p,q)=\{(1-t)p+tq:0<t<1\}.
\]
Under the standard point-visibility convention, the visibility graph of a finite set \(P\subset\R^2\) has vertex set \(P\), with distinct \(p,q\in P\) adjacent exactly when \((p,q)\cap P=\varnothing\). Points of \(P\) are distinct; collinearities are permitted, and a third point on the open segment blocks visibility.

\textbf{Recognition problem (source Problem 2).} Given \(G\), determine whether there exist a finite set \(P\subset\R^2\) and a bijection \(f:V\to P\) such that, for every pair of distinct vertices \(u,v\in V\),
\[
\{u,v\}\in E
\quad\Longleftrightarrow\quad
(f(u),f(v))\cap P=\varnothing.
\]
\textbf{Characterization problem.} Give necessary and sufficient conditions on a finite simple undirected graph for such a representation to exist.

\textit{Qualification.} The supplied excerpt does not explicitly define point visibility or the input graph conventions. The formulation above uses the standard unobstructed-open-segment convention and finite simple graphs; these conventions require confirmation against the cited source. No general-position hypothesis or computational complexity bound is supplied.
\subsection{Short English statement}
Determine whether a graph given by its adjacency matrix can be represented by distinct points in the plane, with an edge exactly when no other representing point lies strictly between its endpoints. Also characterize all graphs admitting such a representation. The registry labels the problem open; that label is supplied metadata, not an independently verified assessment.
\subsection{Sources}
\begin{itemize}
\item Local UnsolvedMath registry, permanent ID \(800533\), exact source problem code \texttt{catalogue entry}, title \emph{Visibility Graphs of Points, Segments and Polygons --- Problem 2}. The supplied record attributes the problem to Ghosh and Goswami, with proposed year 2010, and labels it \texttt{open}.
\item Ghosh and Goswami, \emph{Unsolved Problems in Visibility Graphs of Points, Segments and Polygons} (2010), Problem 2. Supplied source URL: \url{https://arxiv.org/abs/1012.5187}. The background records access on 2026-10-16 and extraction by \texttt{arxiv\_latex}; these provenance details are supplied metadata and were not independently verified here.
\end{itemize}
\end{document}
